\documentclass{article}
\usepackage{enumitem}
\usepackage[margin=1in]{geometry}
\usepackage{titling}
\usepackage{amsmath}
\usepackage{titlesec}

\setlength{\droptitle}{-2cm}
\setlength{\parindent}{0pt}
\setlength{\parskip}{4pt}


\titlespacing*{\section}{0pt}{12pt}{4pt}

\title{
    \textbf{\fontsize{16}{19.2}\selectfont STA380 Project Proposal}\\[0.5em]
    \fontsize{11}{13.2}\selectfont High-Dimensional Feature Selection via Permutation Test and Bootstrap Stability Analysis
}
\author{
    \fontsize{11}{13.2}\selectfont
    Renfei Wu, Shiqi Tian, Chengxi Li, Liwen Yin \\
    \fontsize{11}{13.2}\selectfont
    Team Chou Why Did \\
    \fontsize{11}{13.2}\selectfont
    STA380H5, 2026 \\
    \fontsize{11}{13.2}\selectfont
    University of Toronto Mississauga
}
\date{}

\begin{document}

\maketitle

\section{Project Topic}
This project develops a high-dimensional feature selection method with permutation testing and bootstrap stability analysis. Users will interactively explore how the selected features change as the signal strength, distributional shape, and noise level vary under controlled data-generating processes.

\section{Simulation vs. Dataset}
The project will rely on simulated datasets.

\section{Project Details}
Let $n$ be the sample size and $p$ be the number of features. Let $Y_i \in \{0,1\}$ be the class label for sample $i=1,...,n$. Let $S\subset \{1,...,p\}$ be the set of truly relevant features. 

For each feature $j$ and sample $i$, we generate the data $X_{ij}$ according to a mixture model:
\[
X_{ij} \sim 
\alpha \cdot \mathcal{G}_{Y_i}(\theta_j) + (1-\alpha) \mathcal{H}(\phi_j)
\]
where:
\begin{itemize}[itemsep=0pt, parsep=0pt, topsep=0pt]
    \item $\alpha \in [0,1]$ controls noise strength. When $\alpha=0$, $X_{ij}$ becomes an irrelevant feature.
    \item $\mathcal{G}_{Y_i}(\theta_j)$ is a conditional distribution whose parameters (e.g., mean shift) depend on $Y_i$.
    \item $\mathcal{H}(\phi_j)$ is a noise distribution independent of $Y_i$, so the null features satisfy $X_{\cdot j} \perp Y$.
\end{itemize}

Users will be able to select different distributions for both $\mathcal{G}$ and $\mathcal{H}$ to investigate how data distributions influence results. The ground-truth $S$ and parameters $\theta_j$ will be randomly generated and stored internally for evaluation, while users can only observe the dataset and the selected features.

Feature relevance will be evaluated using permutation testing. For each feature:
\begin{itemize}[itemsep=0pt, parsep=0pt, topsep=0pt]
    \item Users may choose a test statistic $T$ among several test statistics (e.g., mean values).
    \item A null distribution will be constructed by repeatedly shuffling the labels and recomputing $T$ each time.
    \item p-values will be computed based on the permutation distribution:
    \[
    p = \frac{\# (\text{shuffled }T \geq \text{observed }T)}{\#\text{shuffles}}
    \]
\end{itemize}

Bootstrap resampling will be used to quantify the stability of feature selection:
\begin{itemize}[itemsep=0pt, parsep=0pt, topsep=0pt]
    \item Feature selection process is run repeatedly on each bootstrap sample with selection frequency:
    \[
    f_j=\frac{\#\text{times feature $j$ is selected}}{\#\text{bootstrap runs}}
    \]
    \item Features with high selection frequency are more stable.
\end{itemize}

The overall goal is to allow users to explore how permutation tests and bootstrap methods respond to different data structures, including varying signal strength, distribution shapes, and noise levels.

\section{Evaluation}
Denote the ground truth set $S$ and the selected set $\hat{S}$. Since the ground truth is known in the simulation, we evaluate performance using:
\begin{itemize}[itemsep=0pt, parsep=0pt, topsep=0pt]
    \item Empirical precision $= \dfrac{|\hat{S}\cap S|}{|\hat{S}|}$
    \item Empirical recall $= \dfrac{|\hat{S}\cap S|}{|S|}$
\end{itemize}

\section{User Inputs (Shiny Components)}
Users will be able to modify the following:
\begin{enumerate}[itemsep=0pt, parsep=0pt, topsep=0pt]
    \item Choice of distributions for $\mathcal{G}_{Y_i}(\theta_j)$ and $\mathcal{H}(\phi_j)$.
    \item Sample size and number of features.
    \item Selection of permutation test statistic (mean difference, Kolmogorov–Smirnov statistic, or Cramér–von Mises statistic).
    \item Selection of bootstrap descriptive statistics (mean, median, standard error, bias).
    \item Number of bootstrap resamples.
    \item Visualization options and plot parameters.
\end{enumerate}

\nocite{*}
\bibliographystyle{abbrvnat}
\bibliography{bibliography}

\section*{Statement on Generative AI Usage}
We did not use generative AI tools in this proposal. All content, including the mathematical formulation, project design, and writing, was completed independently by the team members.

\end{document}